\section*{अध्याय \ref{ch:b1:elementary-steps} --- अभिक्रिया क्रियाविधियाँ: मूल चरण और सन्निकटन}

\begin{solution}{exo:b1:elementary-steps:1}
(a) एकाणुक: \omterm{def:b1:elementary-steps:elementary-step}{मूल चरण} हो सकता है। (b) द्विअणुक: \omterm{def:b1:elementary-steps:elementary-step}{मूल चरण} हो सकता है।
(c) तीन अणु और एक साथ चार आबंधों का पुनर्निर्माण: \omterm{def:b1:elementary-steps:elementary-step}{मूल चरण} नहीं।
(d) त्रिअणुक, जिसमें \ce{N2} निकली हुई ऊर्जा को बाहर ले जाता है: दुर्लभ पर
वास्तविक \omterm{def:b1:elementary-steps:elementary-step}{मूल चरण}।
\end{solution}

\begin{solution}{exo:b1:elementary-steps:2}
$v = k[\mathrm{A}][\mathrm{B}]$; $v = k[\mathrm{A}]^2$; $v = k[\mathrm{A}]$।
समग्र समीकरण एक लेखा-जोखा है, इसका वर्णन नहीं कि
अणु कैसे मिलते हैं; उसका \omterm{def:b1:rate-laws:order}{दर नियम} मापना पड़ता है, या किसी
क्रियाविधि से व्युत्पन्न करना पड़ता है।
\end{solution}

\begin{solution}{exo:b1:elementary-steps:3}
मध्यवर्ती 30 वाला गर्त है; \omterm{def:b1:elementary-steps:energy-profile}{संक्रमण अवस्थाएँ} 70 और 55 वाले
उच्चिष्ठ हैं। सबसे ऊँची \omterm{def:b1:elementary-steps:energy-profile}{संक्रमण अवस्था} पहली है: पहला
चरण दर-निर्धारक है। यह ऊष्माशोषी है (0 से ऊपर 30 तक)।
\end{solution}

\begin{solution}{exo:b1:elementary-steps:4}
उसी गुणक $10^4$ से (\cref{prop:b1:elementary-steps:catalysis}); साम्य
स्थिरांक अपरिवर्तित रहता है; साम्य लगभग $10^4$
गुना जल्दी प्राप्त होता है।
\end{solution}

\begin{solution}{exo:b1:elementary-steps:5}
$t_{\max} = \ln(0.30/0.10)/(0.30 - 0.10) = \qty{5.49}{s}$। तब
$[\mathrm{B}]_{\max} = a_0\frac{0.10}{0.20}(\eu^{-0.549} - \eu^{-1.648}) =
0.5\,a_0(0.577 - 0.192) = 0.192\,a_0$।
\end{solution}

\begin{solution}{exo:b1:elementary-steps:6}
$\dd[\mathrm{I}]/\dd t = k_1[\mathrm{A}] - (k_{-1} + k_2)[\mathrm{I}] = 0$
से $[\mathrm{I}] = k_1[\mathrm{A}]/(k_{-1} + k_2)$ और $v = k_2[\mathrm{I}]
= \frac{k_1k_2}{k_{-1}+k_2}[\mathrm{A}]$ मिलते हैं। यदि $k_2 \gg k_{-1}$, तो $v =
k_1[\mathrm{A}]$: पहला चरण दर-निर्धारक है। यदि $k_2 \ll k_{-1}$, तो
$v = (k_1/k_{-1})k_2[\mathrm{A}]$: एक पूर्व-साम्य, जिसके बाद एक मंद
चरण आता है।
\end{solution}

\begin{solution}{exo:b1:elementary-steps:7}
योग: \ce{2H2O2 -> 2H2O + O2}। \omterm{def:b1:elementary-steps:catalyst}{उत्प्रेरक}: \ce{I-} (खपता है, फिर
पुनः बनता है); मध्यवर्ती: \ce{IO-}। मंद पहला चरण
दर तय करता है: $v = k_1[\ce{H2O2}][\ce{I-}]$।
\end{solution}

\begin{solution}{exo:b1:elementary-steps:8}
चरण प्रबल रूप से चढ़ाई वाला है: हैमंड अभिगृहीत के अनुसार उसकी \omterm{def:b1:elementary-steps:energy-profile}{संक्रमण
अवस्था} कार्बधनायन जैसी होती है। जो प्रतिस्थापी कार्बधनायन की ऊर्जा घटाता है,
वह \omterm{def:b1:elementary-steps:energy-profile}{संक्रमण अवस्था} की ऊर्जा भी लगभग उतनी ही घटाता है, इसलिए
अवरोध गिरता है और चरण तेज़ होता है।
\end{solution}

\begin{solution}{exo:b1:elementary-steps:9}
$k_{\mathrm B}/k_{\mathrm C} = \exp(20000/RT)$: \qty{300}{K} पर \num{3.0e3},
\qty{600}{K} पर $\exp(4.01) = 55$। गर्म करना B के लिए गतिक वरीयता को घटाता है
और, B के निर्माण को उत्क्रमणीय बनाकर, निकाय को
अधिक स्थायी C तक पहुँचने देता है।
\end{solution}

\begin{solution}{exo:b1:elementary-steps:10}
$k_1[\mathrm{A}][\mathrm{M}] = (k_{-1}[\mathrm{M}] + k_2)[\mathrm{A}^*]$,
इसलिए $v = k_2[\mathrm{A}^*] = \dfrac{k_1k_2[\mathrm{A}][\mathrm{M}]}{k_{-1}[\mathrm{M}]
+ k_2}$। उच्च दाब पर ($k_{-1}[\mathrm{M}] \gg k_2$), $v =
(k_1k_2/k_{-1})[\mathrm{A}]$: कोटि 1। निम्न दाब पर $v =
k_1[\mathrm{A}][\mathrm{M}]$: कोटि 2, टक्कर द्वारा सक्रियण
मंद चरण बन जाता है।
\end{solution}

\begin{solution}{exo:b1:elementary-steps:11}
उपपत्ति प्रमेय वाली ही है। $k_1 = k_2 = k$ के लिए:
$\frac{\dd}{\dd t}([\mathrm{B}]\eu^{kt}) = ka_0$, इसलिए $[\mathrm{B}]\eu^{kt} =
ka_0t$ और $[\mathrm{B}] = a_0kt\,\eu^{-kt}$, जो $t = 1/k$ पर अधिकतम है।
\end{solution}

\begin{solution}{exo:b1:elementary-steps:12}
$\mathrm{A} + \mathrm{B} \rightleftharpoons \mathrm{I} + \mathrm{C}$ (तीव्र,
स्थिरांक $K_1$), फिर $\mathrm{I} + \mathrm{B} \to \mathrm{P}$ (मंद,
$k_2$)। योग समग्र अभिक्रिया है; $[\mathrm{I}] =
K_1[\mathrm{A}][\mathrm{B}]/[\mathrm{C}]$ और $v = k_2[\mathrm{I}][\mathrm{B}]
= k_2K_1[\mathrm{A}][\mathrm{B}]^2/[\mathrm{C}]$: पूर्व-साम्य में मुक्त
उपोत्पाद अभिक्रिया को धीमा करता है।
\end{solution}

\begin{solution}{pb:b1:elementary-steps:1}
\textbf{1.} $[\ce{NO}]_0$ को दोगुना करने से $v_0$, 4 से गुणा होता है: कोटि 2;
$[\ce{O2}]_0$ को दोगुना करने से वह 2 से गुणा होता है: कोटि 1।
\textbf{2.} $v = k[\ce{NO}]^2[\ce{O2}]$; $k = 1.0 \times 10^{-5}/((1.0
\times 10^{-3})^2 \times 1.0 \times 10^{-3}) =
\qty{1.0e4}{L^2.mol^{-2}.s^{-1}}$।
\textbf{3.} तीन अणुओं की एक साथ टक्कर दुर्लभ है, और कोई
अकेला चरण गर्म करने पर धीमा नहीं हो सकता (प्रश्न 5)।
\textbf{4.} $E_a = R\ln(k_{400}/k_{300})/(1/300 - 1/400) = 8.314 \ln(0.1)
/(8.33 \times 10^{-4}) = \qty{-23}{kJ/mol}$।
\textbf{5.} \omterm{def:b1:elementary-steps:elementary-step}{मूल चरण} उन टक्करों से आगे बढ़ता है जो किसी अवरोध को पार करती हैं;
ऐसा कर सकने वाली टक्करों का भिन्न, $\eu^{-E_a/RT}$, जहाँ $E_a \ge
0$, ताप के साथ केवल बढ़ ही सकता है।
\textbf{6.} $-\dd[\ce{NO}]/\dd t = 2v$, $-\dd[\ce{O2}]/\dd t = v$।
\textbf{7.} \ce{2NO -> N2O2} और \ce{N2O2 + O2 -> 2NO2} मिलकर
\ce{2NO + O2 -> 2NO2} देते हैं; मध्यवर्ती \ce{N2O2} है।
\textbf{8.} दोनों द्विअणुक (और पहले चरण का प्रतीप एकाणुक)।
\textbf{9.} $v = k_2[\ce{N2O2}][\ce{O2}]$।
\textbf{10.} $[\ce{N2O2}] = K_1[\ce{NO}]^2$, जहाँ $K_1 = k_1/k_{-1}$।
\textbf{11.} $v = k_2K_1[\ce{NO}]^2[\ce{O2}]$: प्रेक्षित नियम।
\textbf{12.} $k = k_1k_2/k_{-1}$।
\textbf{13.} $\dd[\ce{N2O2}]/\dd t = k_1[\ce{NO}]^2 - k_{-1}[\ce{N2O2}] -
k_2[\ce{N2O2}][\ce{O2}]$।
\textbf{14.} इसे शून्य रखने पर: $[\ce{N2O2}] = k_1[\ce{NO}]^2/(k_{-1} +
k_2[\ce{O2}])$।
\textbf{15.} $v = k_2[\ce{N2O2}][\ce{O2}]$ से दिया गया नियम मिलता है।
\textbf{16.} $k_{-1} \gg k_2[\ce{O2}]$: द्वितय डाइऑक्सीजन से मिलने की तुलना में कहीं
अधिक बार टूट जाता है, इसलिए पहला चरण साम्य में बना रहता है।
\textbf{17.} यदि $k_2[\ce{O2}] \gg k_{-1}$, तो $v = k_1[\ce{NO}]^2$: \ce{NO} में
कोटि 2 और \ce{O2} में 0।
\textbf{18.} \ce{O2} में मापी गई कोटि 1: पहली सीमा।
\textbf{19.} $\ln k = \ln k_1 + \ln k_2 - \ln k_{-1}$।
\textbf{20.} $T$ के सापेक्ष अवकलन करके और $\dd\ln k_i/\dd T
= E_i/RT^2$ का उपयोग करके: $E_a = E_1 + E_2 - E_{-1}$।
\textbf{21.} गर्म करना मंद चरण को थोड़ा तेज़ करता है ($E_2$ छोटी) पर
पूर्व-साम्य को प्रबलता से वापस \ce{NO} की ओर खिसका देता है ($E_{-1}$ बड़ी): कम
द्वितय, और कुल दर घट जाती है।
\textbf{22.} द्वितय अभिकारकों से $E_1 - E_{-1} = \qty{-38}{kJ/mol}$ नीचे है
और दूसरी \omterm{def:b1:elementary-steps:energy-profile}{संक्रमण अवस्था} द्वितय से $E_2 = \qty{15}{kJ/mol}$ ऊपर:
अभिकारकों से \qty{23}{kJ/mol} नीचे। मार्ग का सबसे ऊँचा बिंदु
पहली \omterm{def:b1:elementary-steps:energy-profile}{संक्रमण अवस्था} है, केवल \qty{2}{kJ/mol} ऊपर।
\textbf{23.} $\exp\big(\frac{23000}{8.314}(\frac{1}{400} - \frac{1}{300})\big)
= 0.10$: \qty{400}{K} पर दस गुना धीमी।
\textbf{24.} $E_a = 2 - 40 + 15 = \textbf{\qty{-23}{kJ/mol}}$, वही मान
जो भाग I में मापा गया: क्रियाविधि \omterm{def:b1:rate-laws:order}{दर नियम} और उसकी
विचित्र ताप-निर्भरता, दोनों की व्याख्या करती है।
\end{solution}
